\section{Visual Storytelling}
\label{sec:theory:storytelling}

The visualisation primitives of
Section~\ref{sec:theory:primitives} provide the technical
building blocks of a dashboard. Whether a chart built from
those blocks actually communicates with a stakeholder, however,
depends on a different body of practice --- visual
storytelling. Knaflic's \emph{Storytelling with
Data}~\cite{knaflic2015} is the foundational practitioner
reference; this section summarises the four principles from
that book that govern the visual design of analytical
dashboards.

\subsection{Choose the right chart for the question}

Knaflic argues that the analyst's first decision is the
\emph{type} of chart, and that the canonical chart type for
each analytical question is almost always the right choice.
Time series should be rendered as line or stacked-area charts,
distributions as histograms or box plots, two-variable
relationships as scatter plots, hierarchical part-of-whole as
treemaps, and categorical comparisons as horizontal bars. The
implication is that ornate alternatives --- 3D bars, gauges,
radial layouts --- whose decorative weight outpaces their
information content should be avoided. The analyst's attention
should be on the data and the analytical claim, not on the
novelty of the chart.

\subsection{Eliminate clutter}

Knaflic's second principle is that every visual element on a
chart competes for the analyst's attention; elements that do
not carry information should be removed. Concrete examples
include gridlines that double as axis tick-marks, axis labels
that duplicate what the chart title already says, legends
that duplicate already-visible labels, and decorative drop
shadows or gradients. The principle generalises to dashboard
chrome: filter widgets that show their default values, status
bars that are always green, and breadcrumb navigation that
never changes. Removing such elements is one of the cheapest
visual improvements a dashboard can make, because every
removal increases the relative weight of the remaining
elements.

\subsection{Direct attention}

Knaflic's third principle is that the analyst's eye should be
directed to the analytical claim of the chart, not to its
decorative chrome. Three techniques are most often used.
\emph{Pre-attentive attributes} --- colour, size, position ---
are reserved for the variables being foregrounded; everything
else uses neutral grey. A consistent \emph{colour role} is
maintained across charts: a single primary colour is reserved
for the variable being highlighted, and a diverging palette
(typically red--amber--green) is reserved for ``low / medium /
high'' scales. \emph{Annotation} draws the eye to specific
points: a single labelled callout on a long line chart is
worth more than a dense legend with the same information. The
common thread is that the analyst is told what to look at,
not asked to find it.

\subsection{Tell a story across the dashboard}

The fourth principle, and the one most specific to multi-view
dashboards, is that the analyst's progression from view to
view should follow a narrative. The interaction model
(Section~\ref{sec:theory:mcv}) is the mechanism, but the
narrative is a design choice. The canonical narrative is
\emph{aggregate $\rightarrow$ detail}: the analyst opens on a
high-level KPI, identifies an outlier, clicks to drill into
the outlying entity, inspects its components, and either
resolves the question or escalates back to the aggregate to
compare the entity against its peers. Knaflic frames this as
``drill-from-symptom-to-cause'' and argues that any dashboard
worth using must support it explicitly through one-click
navigation between views, not just through a side-by-side
layout that leaves the navigation up to the analyst.

The chart-design invariant of this thesis, denoted~A2 in
Section~\ref{sec:proposed:invariants}, applies these principles
verbatim: every chart in the dashboard carries a descriptive
title in neutral language and a one-sentence story subtitle in
emphasised type that tells the viewer what to notice in the
chart; default chart colours are kept muted, with saturated
colour reserved for the single element the story subtitle points
at; and chart annotations underline the same finding the subtitle
states.
